\section{Related Work}

\subsection{Model Collapse in Iterative Training}

\citet{shumailov2024ai} demonstrate that language models trained recursively on their own generated data undergo irreversible tail distribution erosion, establishing model collapse as a fundamental limitation of self-consuming training loops. \citet{dohmatob2024tale} further analyze this through scaling laws, showing that synthetic data truncates heavy-tailed input distributions, causing loss of scaling and skill degradation---directly relevant to the rare n-gram tail erosion we observe. \citet{alemohammad2024self} formalize this as Model Autophagy Disorder (MAD), proving theoretically that without fresh real data, iterative self-consumption degrades either quality or diversity. Subsequent work has extended these findings along several axes: \citet{dohmatob2025strong} establish a sharp phase transition in collapse behavior beyond a critical iteration threshold; \citet{seddik2024how} derive the maximum synthetic data proportion before collapse becomes inevitable; and \citet{guo2024curious} track the decay of lexical, syntactic, and semantic diversity across self-training generations, finding the strongest degradation in high-entropy generation tasks.

On the mitigation side, \citet{gerstgrasser2024model} show that accumulating data across generations---rather than replacing it---can prevent collapse, while \citet{briesch2023large} confirm that replacement-based training degrades output diversity proportionally to synthetic data fraction. Our work operates within the replacement setting where collapse is expected, but reveals that a natural mitigation strategy---diversity-based data selection---can paradoxically \emph{accelerate} rather than prevent collapse when it relies on the training model's own representations.

It is instructive to contrast self-consuming language model training with self-play in reinforcement learning \citep{silver2017mastering} and self-training with verification \citep{singh2024beyond}. In both settings, models train iteratively on their own outputs, yet collapse does not occur---because an external reward signal (game outcome, unit test, verifier) anchors the training distribution. Pure language model self-training lacks such an anchor: the model is both generator and (implicit) evaluator, and embedding-based selection provides a proxy signal that degrades endogenously rather than an external ground truth that remains stable.

\subsection{Data Curation and Selection for LLM Training}

Data selection has emerged as a key lever for training efficiency and quality. \citet{sorscher2022beyond} show that metric-guided data pruning can break neural scaling law power laws, motivating a family of embedding-based and model-aware selection methods. SemDeDup \citep{abbas2023semdedup} removes semantic near-duplicates using pretrained embeddings; D4 \citep{tirumala2023d4} combines deduplication with diversity-maximizing selection in embedding space; and DSIR \citep{xie2023data} selects data via importance resampling to match target domain feature distributions. More recent approaches optimize for task-specific signals: LESS \citep{xia2024less} selects instruction-tuning data based on gradient similarity to target tasks; DsDm \citep{engstrom2024dsdm} frames selection as datamodel-based optimization, noting that naive quality proxies can harm downstream performance; QuRating \citep{wettig2024qurating} uses LLM-generated quality scores; and SelectIT \citep{liu2024selectit} leverages model uncertainty for instruction data filtering.

These methods are designed for one-shot curation of static datasets, where the embedding model is fixed and its representation quality is not in question. We show that the iterative self-training setting introduces a qualitatively different failure mode: when the model providing embeddings co-degrades with the training distribution, embedding-based selection amplifies rather than counteracts distributional collapse. Data mixing strategies that blend real and synthetic data across generations \citep{gerstgrasser2024model} offer a complementary mitigation but do not address the selection mechanism itself.

\subsection{Embedding-Based Diversity and Representation Learning}

Submodular maximization provides the algorithmic foundation for embedding-based diversity selection. \citet{iyer2013submodular} establish the optimization framework for submodular functions under budget constraints; \citet{wei2015submodularity} formalize data subset selection via facility location and graph-cut objectives in feature space; and \citet{mirzasoleiman2020coresets} apply gradient-embedding facility location for coreset construction. These methods assume a fixed, high-quality embedding space---an assumption we show breaks down under iterative self-training.

Active learning methods face a related challenge: acquisition functions must remain well-calibrated as the model updates. BADGE \citep{ash2020deep} selects diverse batches via gradient embeddings; BatchBALD \citep{kirsch2019batchbald} maximizes joint mutual information; and subsequent work has shown that standard acquisition strategies degrade under distribution shift \citep{kothawade2021similar} and that embedding quality determines selection effectiveness \citep{yehuda2022active}. However, active learning typically involves a few rounds of model updating, not the sustained multi-generation feedback loop where selection bias compounds into the selecting model's own training data.

The representation learning literature documents several mechanisms by which embedding spaces can lose expressiveness. \citet{jing2022understanding} characterize dimensional collapse in contrastive learning, where embeddings are driven into low-dimensional subspaces; \citet{hua2021feature} identify correlated feature axes as a driver of this collapse; and \citet{papyan2020prevalence} show that deep network representations converge to low-dimensional simplex structures during prolonged training. \citet{laurent2024feature} argue that feature collapse can be benign in classification, but our findings demonstrate that in iterative self-training, the same geometric contraction becomes pathological---it causes diversity selection to invert its intended effect.

\subsection{Representation Drift and Proxy Misalignment}

Our work connects two established literatures that, together, explain why internal-embedding diversity selection fails under iterative training.

The concept drift literature studies how learned representations degrade under sustained distribution shift. \citet{lu2019learning} and \citet{gama2014survey} provide comprehensive taxonomies of drift detection and adaptation. In the continual learning setting, \citet{caccia2022new} quantify abrupt representation change during sequential training, and \citet{davari2022probing} demonstrate that representation-level forgetting proceeds independently of task performance---paralleling our finding that embedding geometry diverges from content diversity even as standard representation quality metrics (RankMe, $k$-NN accuracy) remain stable.

The Goodhart's law and reward hacking literature formalizes the failure of proxy optimization under distribution shift. \citet{gao2023scaling} derive scaling laws showing that overoptimizing a proxy reward model degrades true objective performance; \citet{skalse2022defining} prove that non-trivial proxy rewards are inherently exploitable; \citet{pan2022effects} map how proxy-objective misalignment compounds under increasing optimization pressure; and \citet{karwowski2024goodhart} provide geometric characterizations of when proxy optimization diverges from the true objective. The embedding alignment gap is an instance of this phenomenon: internal embedding diversity serves as a proxy for content diversity, and this proxy progressively loses fidelity as the model's representations co-evolve with the selected data distribution.

Finally, measuring diversity itself requires careful metric design. Distinct-N \citep{li2016diversity} and Self-BLEU \citep{zhu2018texygen} quantify surface-level lexical diversity, while \citet{tevet2021evaluating} audit failure modes of both n-gram and neural diversity metrics. MAUVE \citep{pillutla2021mauve} measures distributional gaps between generated and human text via embedding-space divergence frontiers. We use rare n-gram density as our content diversity measure precisely because, unlike embedding-based metrics, it is model-independent and directly captures the lexical tail erosion that defines model collapse.

\paragraph{Positioning.} The embedding alignment gap as a concept---that proxy metrics degrade under optimization---is well established across the literatures surveyed above. Our contribution is not conceptual novelty but empirical quantification: we provide a controlled factorial ablation (four methods $\times$ three seeds) that decomposes the failure into its constituent factors, a diagnostic metric (selection ratio) that tracks proxy-target divergence in real time, and evidence that the failure is specific to internal embeddings rather than to diversity selection in general.
